\documentclass[11pt,a4paper]{article}

\usepackage[utf8]{inputenc}
\usepackage[T1]{fontenc}
\usepackage[brazil]{babel}
\usepackage{amsmath,amssymb}
\usepackage{booktabs}
\usepackage{geometry}
\usepackage{hyperref}
\usepackage{xcolor}

\geometry{margin=2.5cm}
\hypersetup{colorlinks=true, linkcolor=blue, citecolor=blue, urlcolor=blue}

\title{Hipóteses nulas dos testes estatísticos bidimensionais de energia, Wasserstein e Bhattacharyya}
\author{}
\date{}

\begin{document}

\maketitle

\section{Contexto}

Nas comparações realizadas entre a amostra observada (\texttt{test}) e a
amostra sintética gerada pelo modelo (\texttt{df\_amostras}) no espaço
bidimensional do diagrama BPT (\([\mathrm{N\textsc{ii}}]/\mathrm{H}\alpha\)
vs.\ \([\mathrm{O\textsc{iii}}]/\mathrm{H}\beta\)), utilizam-se três
abordagens para quantificar a similaridade entre as duas distribuições:
o teste de energia de Székely e Rizzo, a distância de Wasserstein e o
coeficiente de Bhattacharyya. Este documento descreve, de forma didática,
a hipótese nula (ou a ausência dela) em cada caso.

\section{Teste de energia (Székely \& Rizzo)}

\textbf{Hipótese nula.} As duas amostras vêm da mesma
distribuição de probabilidade, isto é:
\[
H_0: F_{\text{test}} = F_{\text{amostras}},
\]
onde $F$ denota a função de distribuição conjunta no plano BPT.

\textbf{Intuição.} A estatística de energia mede uma espécie de
``distância elétrica'' entre duas nuvens de pontos, comparando as
distâncias médias \emph{dentro} de cada grupo com as distâncias médias
\emph{entre} os grupos \cite{szekely2004,szekely2013}. Quando as duas
distribuições coincidem, os pontos de um grupo estão, em média, tão
próximos dos pontos do outro grupo quanto dos seus próprios, e a
estatística de energia tende a zero.

\textbf{Inferência.} O p-valor é obtido por permutação: os rótulos de
grupo (\texttt{test} / \texttt{amostras}) são embaralhados repetidamente,
a estatística de energia é recalculada a cada permutação, e a energia
observada é comparada à distribuição empírica resultante. Um p-valor
baixo leva à rejeição de $H_0$, indicando que as distribuições diferem
de forma estatisticamente significativa; um p-valor alto é consistente
com a hipótese de que o modelo reproduz bem a distribuição observada.

\section{Distância de Wasserstein}

Formalmente, a distância de Wasserstein não é um teste de hipótese, mas
uma métrica de transporte ótimo entre duas distribuições
\cite{villani2009,rubner2000}. No entanto, ao ser combinada com um
esquema de permutação, ela passa a admitir uma hipótese nula análoga à
do teste de energia.

\textbf{Hipótese nula.} Assumindo o procedimento de permutação,
\[
H_0: F_{\text{test}} = F_{\text{amostras}},
\]
de modo que qualquer diferença observada na distância de Wasserstein
seria explicável apenas por flutuação amostral.

\textbf{Intuição.} A distância de Wasserstein corresponde ao custo
mínimo de trabalho (massa de probabilidade multiplicada pela distância
percorrida) necessário para transformar uma distribuição na outra
\cite{kantorovich1942,villani2009}. Diferentemente do teste de energia,
ela é sensível não apenas à existência de diferenças, mas também à
magnitude e à direção geométrica do deslocamento de massa entre as duas
nuvens de pontos.

\textbf{Inferência.} Um valor de $0$ indica distribuições idênticas. O
p-valor de permutação é calculado com a mesma lógica do teste de
energia: embaralham-se os rótulos de grupo, recalcula-se a distância de
Wasserstein em cada permutação, e compara-se com o valor observado.

\section{Coeficiente e distância de Bhattacharyya}

\textbf{Ausência de hipótese nula formal.} Ao contrário dos dois testes
anteriores, o coeficiente de Bhattacharyya não constitui um teste de
hipótese estatística com p-valor associado \cite{bhattacharyya1943,kailath1967}.
Trata-se de uma medida descritiva de \emph{sobreposição} entre duas
densidades de probabilidade estimadas via KDE (\textit{kernel density
estimation}) bidimensional, sem uma distribuição nula ou procedimento de
permutação subjacente.

\textbf{Definição.} Dadas duas densidades estimadas $p(x,y)$ e $q(x,y)$
no plano BPT, o coeficiente de Bhattacharyya é definido como
\[
BC = \int\!\!\int \sqrt{p(x,y)\, q(x,y)}\; dx\, dy,
\]
e a distância correspondente é
\[
D_B = -\ln(BC).
\]

\textbf{Interpretação.}
\begin{itemize}
  \item $BC = 1$ ($D_B = 0$): sobreposição total entre as densidades.
  \item $BC = 0$ ($D_B \to \infty$): nenhuma sobreposição entre as densidades.
\end{itemize}

Assim, o coeficiente de Bhattacharyya responde à pergunta ``o quanto as
duas densidades se sobrepõem?'', enquanto os testes de energia e de
Wasserstein respondem, respectivamente, a ``as duas amostras podem ser
consideradas estatisticamente indistinguíveis?'' e ``qual o custo de
transporte entre as duas distribuições, e esse custo é maior do que o
esperado ao acaso?''.

\section{Resumo comparativo}

\begin{table}[h!]
\centering
\begin{tabular}{@{}p{2.7cm}p{5.5cm}p{3cm}p{2.2cm}@{}}
\toprule
\textbf{Método} & \textbf{Pergunta respondida} & \textbf{Tem p-valor?} & \textbf{Referência} \\
\midrule
Teste de energia & As nuvens de pontos são estatisticamente indistinguíveis? & Sim (permutação) & \cite{szekely2004,szekely2013} \\
Wasserstein & Qual o custo de deslocamento de massa entre as distribuições, e é maior que o esperado ao acaso? & Sim (permutação) & \cite{villani2009,kantorovich1942} \\
Bhattacharyya & Quanto as densidades (KDE) se sobrepõem? & Não (medida descritiva) & \cite{bhattacharyya1943} \\
\bottomrule
\end{tabular}
\caption{Comparação entre os três métodos de comparação de distribuições
bidimensionais no diagrama BPT.}
\end{table}

\newpage
\begin{thebibliography}{9}

\bibitem{szekely2004}
Székely, G.\ J., \& Rizzo, M.\ L.\ (2004).
\emph{Testing for equal distributions in high dimension.}
InterStat, 5(16.10), 1--6.

\bibitem{szekely2013}
Székely, G.\ J., \& Rizzo, M.\ L.\ (2013).
\emph{Energy statistics: A class of statistics based on distances.}
Journal of Statistical Planning and Inference, 143(8), 1249--1272.
\url{https://doi.org/10.1016/j.jspi.2013.03.018}

\bibitem{villani2009}
Villani, C.\ (2009).
\emph{Optimal Transport: Old and New.}
Grundlehren der mathematischen Wissenschaften, vol.\ 338.
Springer-Verlag, Berlin Heidelberg.
\url{https://doi.org/10.1007/978-3-540-71050-9}

\bibitem{kantorovich1942}
Kantorovich, L.\ V.\ (1942).
\emph{On the translocation of masses.}
Doklady Akademii Nauk SSSR, 37(7--8), 227--229.

\bibitem{rubner2000}
Rubner, Y., Tomasi, C., \& Guibas, L.\ J.\ (2000).
\emph{The Earth Mover's Distance as a metric for image retrieval.}
International Journal of Computer Vision, 40(2), 99--121.
\url{https://doi.org/10.1023/A:1026543900054}

\bibitem{bhattacharyya1943}
Bhattacharyya, A.\ (1943).
\emph{On a measure of divergence between two statistical populations
defined by their probability distributions.}
Bulletin of the Calcutta Mathematical Society, 35, 99--109.

\bibitem{kailath1967}
Kailath, T.\ (1967).
\emph{The divergence and Bhattacharyya distance measures in signal
selection.}
IEEE Transactions on Communication Technology, 15(1), 52--60.
\url{https://doi.org/10.1109/TCOM.1967.1089532}

\end{thebibliography}

\end{document}
